11269: \subsection{Complementary testing angles}
11270: \label{component:p8f8780350b13}
11271: 
11272: \label{p8f8780350b13:sec:r2-rederived-outer-join}
11273: 
11274: This proof treats every independent testing pair outside the positive
11275: double-infinitesimal angle chart of the original R2 source.  It uses an
11276: explicit full-parent Bell reflection and the two-current and six-root
11277: identities that were freshly proved for R1.  Only those universal algebraic
11278: identities and their real matrix arguments are reused.  No assertion that
11279: the reflected source is an original R1 source is made.
11280: 
11281: \subsubsection*{Actual source and exact reflected parent}
11282: 
11283: The source is the original weak-C2 tuple, with $d=1$, $\lambda=\beta=0$,
11284: the six coefficients in (RER1), all original strict gaps, $n\geq0$,
11285: $a,B,t_s>0$, and $0<r,h<1$.  R2 means
11286: \[
11287: t_s\text{ unbounded},\qquad B, z=aB, y=t_sB\text{ infinitesimal}.
11288: \tag{R2O1}
11289: \]
11290: Put $S=1/t_s$, $l=1-h$, $e=1-r$, $R_s=1+hz$,
11291: $g_s=R_s-ae>0$, and $k=(1+ar)/(1+a)$.  The rederived exact source
11292: lemma (RER1)--(RER7) gives
11293: \[
11294: k\text{ a positive bounded unit},\qquad k<1,\qquad
11295: SC l\sim SAh\sim D/S\sim k/2,
11296: \qquad \Delta/(1+a)\sim1.
11297: \tag{R2O2}
11298: \]
11299: It also gives $ay$ infinitesimal, $Q\sim1$, and the exact expressions
11300: \[
11301: n=\frac{eQ-BK_s}{2r\Delta},\qquad
11302: m=\frac{a^2e(1+y)c s_c}{2g_s\Delta},\qquad
11303: BK_s=(1+a)y(1+hS)(1+z),
11304: \tag{R2O3}
11305: \]
11306: where $c=1+zr\sim1$ and $s_c=1+B(l+ar)\sim1$.
11307: There is no positive lower bound assumed on $r,h,l$, or $g_s$.
11308: 
11309: For arbitrary finite testing angles define the unnormalized testing
11310: state transformation
11311: \[
11312: \begin{split}
11313: G\psi=(&uvW-uX-vY+Z,\ -vW+X-uvY+uZ,\\
11314:        &-uW-uvX+Y+vZ,\ W+vX+uY+uvZ).
11315: \end{split}
11316: \tag{R2O4}
11317: \]
11318: It has $G^TG=D_uD_v I$.  The Bell transformation on observables is
11319: \[
11320: (A'_0,A'_1,B'_0,B'_1)=(-B_1,-B_0,-A_1,-A_0).
11321: \]
11322: After the displayed real local basis gauges it sends the angles to
11323: $(u',v')=(v,u)$ and the source coefficients to
11324: \[
11325: (\lambda',n',m',C',A',D',\alpha',\beta')
11326:  =(\alpha,m,n,C,D,A,\lambda,\beta).
11327: \tag{R2O5}
11328: \]
11329: Direct substitution in the whole parent gives
11330: \[
11331: F_{c'}(v,u;G\psi)=D_uD_v F_c(u,v;\psi).
11332: \tag{R2O6}
11333: \]
11334:   The transformed
11335: $\lambda'=\alpha$ event is kept as its whole positive square.
11336: Because $D_uD_v>0$ is pure testing, a certificate on the left restores
11337: the right by dividing each root by $\sqrt{D_uD_v}$ within $K_t$.
11338: No source or testing contact is identified in this operation.
11339: 
11340: In the next two subsections write the transformed testing angles again
11341: as $u,v$, and denote the reflected coefficients by
11342: \[
11343: \mathsf C=C,\quad\mathsf D=A,\quad\mathsf A=D,
11344: \quad\mathsf m=n,\quad\mathsf n=m.
11345: \]
11346: Thus the symbols on the right refer to the original R2 source.
11347: Define the pure-source quantities
11348: \[
11349: g_* =\frac{\mathsf C\mathsf D}{\mathsf C+\mathsf D},\qquad
11350: h_* =\frac{\mathsf D}{\mathsf C+\mathsf D},\qquad
